\section{From completion to complex linear monodromy}\label{sec:linear}

We prove Theorem~\ref{thm:linear} by applying
Theorem~\ref{thm:completion} on a finite \'etale cover. The inputs from
Cadorel--Deng--Yamanoi concern the total space of the quasi-Albanese map;
they do not require specialness of its general fibre.

We use the following three results of \cite{CDY2025}.
\begin{enumerate}[label=\textup{(\roman*)},leftmargin=2em]
\item By Theorem A, equivalently Theorem 4.1, if $V$ is smooth special
quasi-projective and $\rho:\pi_1(V)\to\GL_d(\C)$ is a representation,
the identity component of the Zariski closure of its image is a direct
product $U\times T$, with $U$ unipotent and $T$ a torus.
\item By Lemma 4.4 and Proposition 4.13, the algebraic quasi-Albanese
map of a smooth special quasi-projective variety is dominant, has
connected general fibre, and has the ordinary exact sequence
\eqref{eq:exact}. The general fibre can be chosen smooth by generic
smoothness. Proposition 4.13 applies to algebraic fibre spaces to
semiabelian varieties whose total source is special (or h-special);
an algebraic fibre space here means a dominant morphism with connected
general fibre. No properness or global surjectivity is required.
\item Connected finite \'etale covers preserve specialness
\cite[Theorem 3.1]{CDY2025}, which attributes this property to
\cite[Proposition 10.11]{Campana2011}.
\end{enumerate}
The specialness assumption of these results is the open-variety
condition in \cite[Definition 2.1]{CDY2025}. The upper-model pair
condition used in Theorem~\ref{thm:example} implies this condition;
we explain the comparison after Definition~\ref{def:special}.

\begin{proof}[Proof of Theorem~\ref{thm:linear}]
Let $L$ be the Zariski closure of $\rho(\pi_1(V))$ and let $L^0$ be its
identity component. The subgroup $\Gamma_0=\rho^{-1}(L^0)$ has finite
index in $\pi_1(V)$. By the algebraic form of the Riemann existence
theorem~\cite[Expos\'e~XII, Th\'eor\`eme~5.1, p.~251]{SGA1}
it corresponds to a connected finite \'etale cover $V_0\to V$.
By the first input, $L^0\simeq U\times T$. Project the restricted
representation of $\pi_1(V_0)$ to $U$.

The third input makes $V_0$ special. Apply the second input to the
quasi-Albanese map of $V_0$ itself: it is dominant, its general fibre
is smooth and connected, and its ordinary fundamental-group sequence
is exact. All the hypotheses of Theorem~\ref{thm:completion} therefore
hold for $V_0$. The projected unipotent image has class at most two,
while the projection to $T$ is abelian. Every triple commutator of
$\rho(\Gamma_0)$ is consequently trivial in both factors of
$U\times T$, and is itself trivial. Since $\rho(\Gamma_0)$ has finite
index in $\rho(\pi_1(V))$, this proves the theorem.
\end{proof}

Reapplying exactness on $V_0$ is essential: unipotent completion need
not be unchanged on passing to a finite-index subgroup. This proof
uses neither fibre-specialness nor a mixed Hodge structure on an
arbitrary linear representation quotient. The mixed Hodge argument
was carried out entirely on the canonical quotients in
Section~\ref{sec:mhs}.

Applied to $V$ itself, the second CDY input also gives the completion
conclusions of Theorem~\ref{thm:completion} for every smooth special
quasi-projective variety. We record the resulting specialization to logarithmic Kodaira
dimension zero. Its linear-image assertion is a specialization of the
theorem of Cao--Deng--Hacon--P\u aun \cite[Theorem A]{CDHP2026}.

\begin{corollary}[Log-Kodaira-zero monodromy]\label{cor:log-kodaira-zero}
Let $V$ be a connected smooth complex quasi-projective variety with
logarithmic Kodaira dimension $\bar\kappa(V)=0$, and put
$\Gamma=\pi_1(V)$. For every positive integer $d$ and every
representation $\rho:\Gamma\to\GL_d(\C)$, the image $\rho(\Gamma)$ has
a subgroup of finite index and nilpotency class at most two. The full
rational unipotent completion of $\Gamma$ is a finite-dimensional
unipotent algebraic group of class at most two. Every finite-dimensional
complex unipotent representation of $\Gamma$ also has image of class
at most two.
\end{corollary}

\begin{proof}
Choose a smooth projective compactification $Y$ of $V$ with reduced
simple normal crossing boundary $D=Y\setminus V$. Then
$\bar\kappa(V)=\kappa(Y,K_Y+D)=0$. Campana's theorem
\cite[arXiv version, Theorem 6.7 and Corollary 4.11]{Campana2011}
applies to smooth projective pairs with reduced simple normal crossing
boundary and gives
specialness of $(Y,D)$. By the comparison after
Definition~\ref{def:special}, $V$ is therefore special in the
open-variety convention of Cadorel--Deng--Yamanoi.

Theorem~\ref{thm:linear} gives the linear-image assertion. Applied to
$V$ itself, the second CDY input above gives a dominant quasi-Albanese
map with smooth connected general fibre and the full ordinary exact
sequence~\eqref{eq:exact}. Theorem~\ref{thm:completion} therefore gives
the assertions about the full rational completion and complex unipotent
images.
\end{proof}

The finite-index subgroup in Corollary~\ref{cor:log-kodaira-zero} may
depend on the representation. The corollary does not assert virtual
two-step nilpotence of the ordinary discrete group $\pi_1(V)$ itself;
the faithfulness issue remains as explained in Section~\ref{sec:ordinary}.
